% 图3: ECMP vs 包喷洒对比示意图
\begin{tikzpicture}[scale=0.9, transform shape,
    box/.style={rectangle, draw, minimum width=1.5cm, minimum height=0.8cm, align=center, font=\small},
    flow/.style={rectangle, draw, minimum width=0.6cm, minimum height=0.4cm, font=\tiny},
    arrow/.style={->, >=stealth, thick},
    label/.style={font=\small\bfseries},
    link/.style={line width=2pt},
    conflict/.style={draw=red, line width=3pt, dashed}
]

% === 左侧: ECMP ===
\node[label] at (-4, 6) {\large \textbf{ECMP: 基于流的负载均衡}};
\node[font=\scriptsize, text width=4cm, align=center] at (-4, 5.5) {相同五元组的数据包走同一条路径\\Hash(源IP, 目的IP, 端口, 协议) → 路径};

% 源节点
\node[box, fill=blue!20] (src1) at (-7, 4) {流A};
\node[box, fill=blue!20] (src2) at (-4, 4) {流B};
\node[box, fill=blue!20] (src3) at (-1, 4) {流C};

% 中间交换机
\node[box, fill=gray!30, minimum width=2cm] (sw-ecmp) at (-4, 2.5) {ECMP交换机\\Hash选路};

% 多条路径
\draw[link, green!50] (-5.5, 1.5) -- (-5.5, 0.5);
\draw[link, yellow!50] (-4, 1.5) -- (-4, 0.5);
\draw[link, orange!50] (-2.5, 1.5) -- (-2.5, 0.5);

% 目标节点
\node[box, fill=green!20] (dst-ecmp) at (-4, -0.5) {目标};

% ECMP数据流
\draw[arrow, blue!60] (src1) -- (-6, 3) -- (-5.5, 3) -- (-5.5, 2.5);
\draw[arrow, blue!60] (src2) -- (-4, 2.5);
\draw[arrow, blue!60] (src3) -- (-2, 3) -- (-2.5, 3) -- (-2.5, 2.5);

% 显示哈希冲突问题
\node[conflict, fit={(-5.5, 1.5) (-5.5, 0.5)}, inner sep=0pt] {};
\node[font=\scriptsize, text=red] at (-6.5, 1) {哈希冲突!};
\node[font=\tiny, text=red] at (-6.5, 0.6) {流A和流X\\走同一路径};

% 流量标注
\node[font=\tiny] at (-5.5, 2.2) {80\%负载};
\node[font=\tiny] at (-4, 2.2) {10\%负载};
\node[font=\tiny] at (-2.5, 2.2) {10\%负载};

% 问题标注
\node[rectangle, draw=red, fill=red!10, text width=3.5cm, align=center, font=\scriptsize] at (-4, -2) {
    \textbf{ECMP问题}\\
    • 大象流哈希冲突\\
    • 负载不均严重\\
    • 链路利用率低\\
    • AI训练性能下降30\%+
};

% === 右侧: 包喷洒 ===
\node[label] at (4, 6) {\large \textbf{包喷洒: 基于包的负载均衡}};
\node[font=\scriptsize, text width=4cm, align=center] at (4, 5.5) {同一流的不同数据包分散到多条路径\\Round-robin / 随机选择 → 路径};

% 源节点（同一个流被分割）
\node[box, fill=purple!20, minimum width=3cm] (src-ps) at (4, 4) {大象流 (分解为包)};

% 中间交换机
\node[box, fill=gray!30, minimum width=2cm] (sw-ps) at (4, 2.5) {包喷洒交换机\\轮询选路};

% 多条路径（均衡负载）
\draw[link, green!70] (2.5, 1.5) -- (2.5, 0.5);
\draw[link, green!70] (4, 1.5) -- (4, 0.5);
\draw[link, green!70] (5.5, 1.5) -- (5.5, 0.5);

% 目标节点 + 重排序缓冲区
\node[box, fill=purple!20] (dst-ps) at (4, -0.5) {目标 + 重排序};

% 包喷洒数据流（轮询）
\draw[arrow, purple!60] (3.2, 3.5) -- (2.5, 3.5) -- (2.5, 2.5);
\draw[arrow, purple!60] (4, 3.5) -- (4, 2.5);
\draw[arrow, purple!60] (4.8, 3.5) -- (5.5, 3.5) -- (5.5, 2.5);

% 包编号标注
\node[font=\tiny, fill=yellow!30] at (2.2, 3.5) {Pkt1};
\node[font=\tiny, fill=yellow!30] at (3.7, 3.5) {Pkt2};
\node[font=\tiny, fill=yellow!30] at (5.8, 3.5) {Pkt3};

% 均衡负载标注
\node[font=\tiny, text=green!60!black] at (2.5, 2.2) {33\%负载};
\node[font=\tiny, text=green!60!black] at (4, 2.2) {33\%负载};
\node[font=\tiny, text=green!60!black] at (5.5, 2.2) {33\%负载};

% 重排序标注
\draw[arrow, orange!60, dashed] (2.5, 0.5) -- (2.8, 0);
\draw[arrow, orange!60, dashed] (4, 0.5) -- (4, 0);
\draw[arrow, orange!60, dashed] (5.5, 0.5) -- (5.2, 0);
\node[font=\tiny, text=orange!70!black] at (6.2, 0.2) {乱序到达};
\node[font=\tiny, text=orange!70!black] at (6.2, -0.3) {硬件重排序};

% 优势标注
\node[rectangle, draw=green, fill=green!10, text width=3.5cm, align=center, font=\scriptsize] at (4, -2) {
    \textbf{包喷洒优势}\\
    • 完美负载均衡\\
    • 链路利用率高(90\%+)\\
    • 适应非对称拓扑\\
    • 智能网卡重排序
};

% === 底部对比总结 ===
\node[rectangle, draw, fill=gray!10, text width=14cm, align=center, font=\small] at (0, -4) {
    \textbf{关键洞察}: AI训练的"低熵"特性（少量大象流）使ECMP哈希冲突成为性能瓶颈。\\
    包喷洒通过细粒度负载均衡充分利用带宽，配合智能网卡解决乱序问题。
};

\end{tikzpicture}
